% Paper blocks for the extraction robustness and prompt ablation (PROJECT.md 9.6).
%
% AuthorKit27/ is git-ignored, so the paper cannot travel through a branch or PR.
% These blocks are version-controlled here and applied by
% `analysis/apply_extraction_snippets.py`, which is idempotent and writes a .bak.
%
% Every number is transcribed from results/extraction_challenge_summary.csv,
% results/extraction_challenge_raw.csv (+ _run1), results/extraction_cost_latency.csv
% and analysis/real_paper_extraction_summary.csv.
%
% Anchors were checked against main.tex at the time of writing: Q1 is the
% real-paper evaluation, Q3 is LLM vs. regex, Q5 is cost/latency, and the new
% material lands as Q6 immediately after the Q5 cost table.


% ==================================================== BLOCK: contribution-bullet
% REPLACES the third contribution bullet in the introduction. The old text
% promised token/cost/latency effects only; the measurement also covers quality,
% and it changed which prompt ships.

    \item \textbf{Cost-aware agentic prompting, measured and revised.}
    Autonomous loops make many LLM calls. We measure a compressed prompting
    scheme (the ``Caveman Protocol'') against a schema-matched natural-language
    prompt on token count, cost, latency \emph{and} extraction quality across
    three commercial providers, and report a negative result: the compression
    costs accuracy for negligible saving, so the deployed system does not use it.


% ==================================================== BLOCK: cost-aware-prompting
% REPLACES the "Cost-Aware Prompting" subsection body in the Approach section.
% The old text asserted that the extraction prompt *uses* the Caveman Protocol
% (no longer true) and cross-referenced Q4 (cost/latency is Q5).

The extraction system prompt was originally written in a compressed, filler-free
style (the ``Caveman Protocol'': no articles, pleasantries, or hedging), applied
only to internal prompts while user-facing output remained in natural language.
We measure its effect rather than assert it (Q5 and Q6). Because that measurement
showed a quality cost that outweighed the saving, the deployed extractor now uses
the natural-language prompt; we keep the compressed variant documented and
measured in Appendix~\ref{app:caveman}.


% ==================================================== BLOCK: q3-pointer
% APPENDED to the end of the Q3 (LLM vs. regex) paragraph. Q3 asserts that the
% LLM's value is robustness on messy prose; Q6 is where that is actually measured.

Q6 tests that claim directly on deliberately messy text, where the gap is stark
rather than absent.


% ==================================================== BLOCK: q5-paragraph
% REPLACES the tail of the Q5 paragraph, from "Google's Gemini" to the end. Fixes
% the superseded latency claim and the 21.9% figure in one edit.

Google's Gemini~2.5~Flash is the cheapest ($\$0.0006$/call) and among the fastest
($2.15$\,s, level with GPT-4o at $2.00$\,s); Anthropic's Claude Sonnet~4.5 is the
most expensive ($\$0.0042$/call) and slowest ($3.95$\,s). We separately measure
the Caveman Protocol's prompt compression at a \textbf{38.6\%} reduction in
system-prompt tokens ($308 \rightarrow 189$), or $30.0\%$ of the prompt tokens
actually billed on a challenge-set request. An earlier version of this
measurement reported $21.9\%$ against a natural-language baseline that omitted
the categorical distribution and its \texttt{categories} field; that baseline
requested strictly less than the Caveman prompt and so \emph{understated} the
compression. Compression does not translate into proportional saving: because
completion tokens dominate this workload, the measured cost reduction is only
$3.5\%$ (paired Wilcoxon $p=.0065$) and latency is unchanged. Q6 shows that it
also costs accuracy.


% ==================================================== BLOCK: cost-table-rows
% REPLACES the three data rows of Table~\ref{tab:cost}. Re-measured on the
% natural-language prompt now shipped (results/extraction_cost_latency.csv, 5 reps).
% The historical Caveman-prompt measurement is preserved in results/cost_latency.csv.

Google    & 2.15 & 0.0006 & 367 & 194 \\
OpenAI    & 2.00 & 0.0023 & 363 & 140 \\
Anthropic & 3.95 & 0.0042 & 379 & 201 \\


% ==================================================== BLOCK: q6-subsection
% INSERTED after the Q5 cost table, before the "Disclosure-risk threat model"
% subsection.

\subsection{Q6: Does Prompt Compression Cost Extraction Quality?}
The compression in Q5 is only meaningful alongside its effect on what is
extracted. We built a 30-case challenge set of deliberately messy source text
spanning ten failure modes: negated findings, several populations reported for
one variable, values that conflict between abstract and table, repeated
timepoints, mixed units, confidence intervals, footnote-marked table cells,
scientific notation, irrelevant nearby numbers, and standard errors presented
where the schema wants a standard deviation. Eleven cases are verbatim excerpts
from the sources of Q1, sliced programmatically at recorded character offsets so
the fixture cannot drift from its sources; nineteen are constructed and labelled
as such. Each case also annotates the values that must \emph{not} be extracted
(p-values, sample sizes, instrument ranges, thresholds), which is what makes
context errors and hallucinations countable rather than invisible. Abstention is
rewarded rather than punished: emitting nothing yields no false positives, so the
three cases whose correct answer is silence score $F_1{=}1$ for an extractor that
stays silent.

The compression is not free (Table~\ref{tab:extraction}). Averaged over three
models and five repetitions, the Caveman prompt reaches $F_1=.710$ against $.778$
for the schema-matched natural-language prompt, a paired deficit of $-.069$
(Wilcoxon $p=.0001$, $n=90$ case\,$\times$\,model pairs) that holds in the same
direction for all three models. The loss is precision, not coverage: recall is
near-ceiling for both ($.960$ vs.\ $.989$) while precision falls from $.720$ to
$.639$, and the hallucination rate rises from $.098$ to $.154$. Stripping the
prose that tells the model what \emph{not} to do makes it invent more numbers.
Against a $3.5\%$ cost saving that is a poor trade, so AgentDataset now ships the
natural-language prompt; we retain the Caveman Protocol as a documented, measured
design because the negative result is itself the contribution---prompt
compression is widely assumed to be free when the schema is held fixed, and on
messy source text it is not.

Both prompts remain far ahead of the regex fallback ($F_1=.156$), whose patterns
require literal ``mean\,\ldots\,std'' phrasing and so rarely fire on real tables;
its precision of $1.000$ reflects near-total silence, which also gives it the best
abstention behaviour of the three arms ($3/3$ silence cases, against $11\%$ for
the natural-language prompt and $0\%$ for the Caveman prompt). Neither LLM arm
reliably declines to answer when a document supports no answer, which we regard
as the most consequential remaining weakness of the extraction stage. Real
excerpts are markedly harder than constructed ones ($F_1$ $.550$ vs.\ $.803$ for
the Caveman arm), the clearest evidence we have that synthetic evaluation text
overstates extraction quality. Appendix~\ref{app:extraction} reports a full
independent replication and the effect of the prompt change on Q1.

\begin{table}[t]
\centering
\setlength{\tabcolsep}{4pt}
\begin{tabular}{lrrr}
\toprule
& Caveman & Natural & Regex \\
\midrule
$F_1$                & .710 & \textbf{.778} & .156 \\
Precision            & .639 & \textbf{.720} & 1.000 \\
Recall               & .960 & \textbf{.989} & .150 \\
Hallucination rate   & .154 & \textbf{.098} & .000 \\
Context-error rate   & .115 & \textbf{.101} & .000 \\
Correct abstention   & .000 & .111 & \textbf{1.000} \\
\midrule
Prompt tokens        & \textbf{283} & 405 & 0 \\
Cost (\$/call)       & \textbf{.0024} & .0025 & 0 \\
\bottomrule
\end{tabular}
\caption{Extraction over the 30-case challenge set, averaged over three
OpenRouter-routed models and five repetitions (930 LLM extractions, 30 regex).
The Caveman prompt buys a $30\%$ prompt-token reduction for $-.069$ $F_1$
(paired Wilcoxon $p=.0001$) and a $57\%$ relative increase in hallucination,
while saving $3.5\%$ of cost. ``Correct abstention'' is measured on the three
cases whose correct answer is to extract nothing.}
\label{tab:extraction}
\end{table}


% ==================================================== BLOCK: appendix-caveman-intro
% REPLACES the two sentences under \label{app:caveman}. The old text described the
% compressed prompt as the extraction prompt and claimed the baseline had
% "identical semantics", which was not true of the baseline then in use.

The compressed extraction system prompt is reproduced below. It is no longer the
deployed prompt: Q6 shows that it costs $-.069$ $F_1$ against the natural-language
prompt for a $3.5\%$ cost saving, so the shipped extractor uses the latter. Both
prompts are included in the released code and request an identical schema, which
the earlier baseline did not---it omitted the categorical distribution and the
\texttt{categories} field, and a regression test now enforces field-set parity
between the two.


% ==================================================== BLOCK: appendix-extraction
% INSERTED after the Caveman verbatim block in the appendix (unlimited length).

\section{Extraction Challenge Set and Replication}
\label{app:extraction}

\paragraph{Protocol.} All arms run at temperature $0$ with
\texttt{response\_format} set to \texttt{json\_object}, the production setting,
and with retries disabled so nondeterminism and transient failure stay visible.
The regex fallback is deterministic, runs once, and is reported as its own arm;
fallback output is never folded into LLM results. Matching is by value rather
than by name, because the regex fallback emits generic identifiers and would
otherwise be penalised for a reason unrelated to reading the numbers correctly.
False positives are separated into context errors (a value present in the
document but annotated as not a distribution parameter) and hallucinations (a
value appearing nowhere in the source text), the latter grounded in an
enumeration of every numeric literal the excerpt contains.

\paragraph{Replication.} We ran the full sweep twice, independently, at identical
settings; both runs are released. The Caveman deficit is $-.069$ in the first and
$-.063$ in the second (Wilcoxon $p=.0001$ and $p=.0004$); pooled arm means agree
to within $.010$ $F_1$ ($.710$ vs.\ $.719$ Caveman, $.778$ vs.\ $.782$
natural-language), and the hallucination gap is present in both. The regex arm is
bit-identical between runs, confirming the observed spread is model-side rather
than harness-side. The second run additionally lost $15\%$ of its Anthropic calls
to transient API errors; because retries are disabled and failures are recorded
as rows rather than dropped, this shrinks the paired sample from 90 to 85 without
biasing the reported rates.

\paragraph{Effect on Q1.} Re-running the real-paper evaluation under the new
prompt on identical cached source text replicates the finding independently:
extraction $F_1$ rises from $.247$ to $.308$ and precision from $.167$ to $.226$,
while recall is unchanged at $.480$. The item counts show the mechanism directly:
false positives fall sharply ($27 \rightarrow 12$ on one source) while true
positives are identical. The natural-language prompt wins by inventing less, not
by finding more. This too was run twice ($F_1$ $.308$ then $.304$). Absolute
performance on real papers remains low, and we do not claim otherwise.

\paragraph{Two measurement caveats.} First, because \texttt{json\_object} is
requested, the malformed-output rate we report counts schema violations rather
than invalid JSON; it is $1.3\%$ for the Caveman arm and $0\%$ for the
natural-language arm. Second, Anthropic models routed through OpenRouter ignore
the \texttt{json\_object} flag and wrap the object in a Markdown code fence on
$100\%$ of successful calls (255/255, against 0/600 for the other two providers).
A strict JSON parse therefore fails on every such response, and a pipeline with a
regex fallback will silently degrade to that fallback for one of three providers
while appearing to succeed. We strip the fence before parsing and report the
incidence rather than hiding it, because the failure is invisible to any
evaluation that does not separate fallback output from LLM output.

\paragraph{The natural-language prompt.} The deployed extraction prompt, which
requests the identical schema to the compressed prompt above:

\begin{quote}\begin{scriptsize}\begin{verbatim}
You are an expert statistician with many years
of experience analyzing research papers and
extracting statistical information from them.
Please carefully read through the following
text and identify all of the statistical
variables that are mentioned, along with their
probability distributions (such as normal,
uniform, gamma, or categorical), their means,
and their standard deviations. Also, please
identify any correlations that are described
between pairs of variables, including the
strength and direction (positive or negative)
of each correlation.

Please format your response as a JSON object
that matches the following schema exactly:
{ "variables": {
    "<name>": {"distribution":
      "normal|uniform|gamma|categorical",
      "mean": 0.0, "std": 1.0, "min": null,
      "max": null, "categories": null} },
  "correlations": {
    "<key>": {"var1": "<name>", "var2": "<name>",
      "correlation": 0.5,
      "direction": "positive|negative"} } }

The "categories" field should only be used when
the distribution is "categorical". In that case
it should be an object that maps each category
label to its probability, for example
{"yes": 0.3, "no": 0.7}.

Please make sure to only output the JSON object
itself and nothing else - no introductory text,
no explanations, and no markdown code fences.
\end{verbatim}\end{scriptsize}\end{quote}
